\section{Separate development calibration}
\label{sec:dev40-results}
The development stage precedes the main sample and uses 40 disjoint randomly selected tasks, the same five conditions and two fresh labelled repeats (2 and 3). All 400 assigned programs and 720 CLI turns complete and pass the trace audit. The full 400 native records reconcile to the exact exported programs. Reference task /1005 fails in the development environment, leaving 390 evaluable attempts over 39 tasks; its ten generation costs remain included. No final-format extraction fails. Each attempt is scored separately, with no best-of-two selection. For paired development contrasts, the two repeats are averaged within task before resampling.

\begin{table}[htbp]\centering\small
\caption{Development results: 80 assigned candidates per condition, two repeats on 40 tasks. Quality uses 78 evaluable attempts; resource means use all 80. Missingness ranges use the full assignment denominator and are not confidence intervals.}\label{tab:dev40-main}
\begin{tabular}{lrrrrrr}\toprule
Arm & Passes & Rate (\%) & Range (\%) & Tokens & USD & No cache\\
\midrule
D & 42/78 & 53.85 & 52.50--55.00 & 4\,500 & 0.00776 & 0.00939\\
SN & 34/78 & 43.59 & 42.50--45.00 & 5\,598 & 0.01251 & 0.01417\\
SR & 39/78 & 50.00 & 48.75--51.25 & 5\,189 & 0.01070 & 0.01234\\
MN & 39/78 & 50.00 & 48.75--51.25 & 15\,636 & 0.02624 & 0.03090\\
MR & 36/78 & 46.15 & 45.00--47.50 & 15\,375 & 0.02543 & 0.03030\\
\bottomrule\end{tabular}\end{table}


The observed rates range from 43.59\% to 53.85\%. D passes 42/78 attempts, SR and MN each 39/78, MR 36/78, and SN 34/78. D also has the smallest resource mean. Complete generation uses 3,703,821 tokens, valued at USD 6.6113451, or USD 7.7687595 without cache discounts. The denominator remains 40 task clusters for resources and 39 for quality contrasts.

\begin{table}[htbp]\centering\small
\caption{Paired development contrasts. Quality differences and descriptive 95\% task-bootstrap intervals are percentage points over 39 complete task clusters; resource ratios use 40. The last two comparisons are exploratory. No confirmatory p-value or non-inferiority test is assigned.}\label{tab:dev40-contrasts}
\begin{tabular}{lrrr}\toprule
Contrast & Quality difference [interval] & Token ratio & USD ratio\\
\midrule
SR - SN & $+6.41\;[+1.28,+11.54]$ & 0.927 & 0.855\\
MR - MN & $-3.85\;[-11.54,+2.56]$ & 0.983 & 0.969\\
MN - SN & $+6.41\;[-2.56,+15.38]$ & 2.793 & 2.097\\
MR - SR & $-3.85\;[-8.97,+1.28]$ & 2.963 & 2.377\\
SN - D & $-10.26\;[-19.23,-2.56]$ & 1.244 & 1.613\\
MR - D & $-7.69\;[-17.95,+1.28]$ & 3.417 & 3.277\\
\bottomrule\end{tabular}\end{table}


The role-minus-neutral development differences are $+6.41$ percentage points in one call and $-3.85$ in three calls. The difference of differences is $-10.26$ points, with a descriptive task-bootstrap interval of $[-19.23,-1.28]$. Three-call token means are 2.79 and 2.96 times the corresponding single-call means; valuation ratios are 2.10 and 2.38. These development observations motivated the separate reserved-task allocation. They are not additional observations in its four-test family and cannot replace a main-sample result. Complete two-repeat outcomes are retained in the appendix and public archive.
